% supplement.tex: Supplementary material of the ICORR 2027 paper. ICORR takes no supplementary upload, so it lives
% here, and the paper points to this repository. Fig. S1 is ../figures/fig4m.tex with the data in ../figures/data/.
% Tables S1-S3 condense CONDITIONS.md; keep both in step. Build: latexmk -pdf supplement.tex
\documentclass[10pt,a4paper]{article}
\usepackage[margin=2.5cm]{geometry}
\usepackage[T1]{fontenc}
\usepackage{newtxtext,newtxmath}
\usepackage{amsmath}
\usepackage{xcolor}
\usepackage{booktabs,tabularx,xltabular}
\usepackage{algorithm}
\usepackage[noend]{algpseudocode}
\usepackage{pgfplots}
\pgfplotsset{compat=1.18}
\usepgfplotslibrary{fillbetween}
\usetikzlibrary{arrows.meta}
\usepackage[labelsep=period,font=small,labelfont=bf]{caption}
\captionsetup[figure]{name=Fig.}
\usepackage{placeins}
\usepackage[hidelinks]{hyperref}
% float pages fill from the top instead of spreading their floats over the page
\makeatletter
\setlength{\@fptop}{0pt}\setlength{\@fpsep}{18pt}\setlength{\@fpbot}{0pt plus 1fil}
\makeatother

% icorr_style.tex: shared colours and line styles of the amplitude-free paper's figures
% (paper/notes/story_concept.md section 7a, scheme C). \input once in the preamble, after
% \usepackage{pgfplots}. Every figure uses these names; no figure defines its own colours.
%
% Roles:
%   PsyPref (all variants) ............ icBlue, solid
%   best published comparator ......... icOrange, solid, hollow marker: square CMA-ES, triangle Arens' GP
%   GP-BALD (same prior) .............. icGrey, dashed
%   CMA-ES on Schaefer's family ....... icDark, dash-dot
%   generic assistance (50 %) ......... icGeneric, solid thin
%   best K-weight profile (ceiling, floor) icContext, dotted
%   K = 8 / 16 / 24 (or k = 24 / 40 / 150) icBlueLight / icBlue / icBlueDark
%   settings as markers ............... circle 3-35 Nm, square fixed 12 Nm, diamond 3-12 Nm
%   robustness: filled = fixed yardstick (post hoc), hollow = registered scoring
\definecolor{icBlue}{HTML}{1845FB}
\definecolor{icBlueLight}{HTML}{8FA4FD}
\definecolor{icBlueDark}{HTML}{0F2C9E}
\definecolor{icOrange}{HTML}{F89C20}
\definecolor{icOrangeText}{HTML}{B86E00}% text in the comparator's hue (contrast 4.0 on white)
\definecolor{icGrey}{HTML}{7A7A7A}
\definecolor{icDark}{HTML}{333333}
\definecolor{icGeneric}{HTML}{9C9CA1}
\definecolor{icContext}{HTML}{BBBBBB}
\definecolor{icRule}{HTML}{DDDDDD}
\definecolor{icInk}{HTML}{222222}
\definecolor{icMuted}{HTML}{666666}
% Fig. 1 only: the pair shown to the user (author's choice, 2026-10-29: distinct hues read
% stronger than two blues; gold was dropped because it equals the comparator's orange)
\definecolor{icTintBlue}{HTML}{F1F3F8}% Fig. 1: the loop, PsyPref's and the user's box (author's choice)
\definecolor{icPairA}{HTML}{BB5566}
\definecolor{icPairB}{HTML}{33BBEE}
\tikzset{
  ic psypref/.style={icBlue, line width=1.5pt, line cap=round, line join=round},
  ic best/.style={icOrange, line width=1.2pt, line cap=round, line join=round},
  ic gpbald/.style={icGrey, line width=0.9pt, dash pattern=on 3.2pt off 1.6pt},
  ic sch/.style={icDark, line width=0.9pt, dash pattern=on 3.6pt off 1.4pt on 0.9pt off 1.4pt},
  ic generic/.style={icGeneric, line width=0.6pt},
  ic context/.style={icContext, line width=0.8pt, dash pattern=on 0.9pt off 1.5pt},
  ic rule/.style={icRule, line width=0.5pt},
  ic label/.style={inner sep=0.6pt, font=\scriptsize, align=left},
  ic note/.style={inner sep=0.8pt, font=\scriptsize, icInk, fill=white, fill opacity=0.85,
                  text opacity=1, align=right},
  ic band/.style={fill=#1, fill opacity=0.15, draw=none},
}
\pgfplotsset{
  ic axis/.style={
    axis line style={icMuted, line width=0.5pt},
    tick style={icMuted, line width=0.5pt},
    tick label style={font=\scriptsize, icMuted},
    label style={font=\footnotesize, icInk},
    title style={font=\footnotesize, icInk},
    axis x line*=bottom, axis y line*=left,
    tick align=outside, major tick length=2pt,
    clip=false,
  },
}

\newcommand{\icorrtikzdata}{../figures/data/}
% fig3p_text.tex: written by scripts/paper_v2_fig3p.py (repository); do not edit by hand.
% Every number of Fig. 3's caption and the draft Results paragraph, formatted from values the script checked
% (fam_report.md, fam_numbers.json, amp_table1.json, Fig. 2's data). Differences as fam_report.md prints them.
\def\pjK{24}
\def\pjNusers{200}
\def\pjNsubj{5}
\def\pjSigma{1}
\def\pjLapse{0.1}
\def\pjGapW{18.8}
\def\pjGapWci{(14.7, 22.8)}
\def\pjGapF{9.0}
\def\pjGapFci{(6.1, 11.9)}
\def\pjSchCeilW{94.1}
\def\pjSchCeilF{88.9}
\def\pjSchPsyW{86.4}
\def\pjSchPsyF{74.8}
\def\pjCmaW{67.6}
\def\pjCmaWci{(63.1, 71.9)}
\def\pjCmaF{65.8}
\def\pjCmaFci{(63.5, 67.9)}
\def\pjPubSch{76.6}
\def\pjPubSchci{(74.6, 78.5)}
\def\pjPubPchip{86.6}
\def\pjPubPchipci{(85.0, 88.3)}
\def\pjPubSchDiff{17.5}
\def\pjPubSchDiffci{(15.6, 19.5)}
\def\pjPubPchipDiff{7.4}
\def\pjPubPchipDiffci{(5.9, 9.0)}
\def\pjProjSchW{6.1}
\def\pjProjSchWci{($-7.8$, $-4.6$)}
\def\pjProjSchWsigned{$-6.1$}
\def\pjProjFranksW{$-0.6$}
\def\pjProjPchipW{$-1.1$}
\def\pjProjFixLo{3.4}
\def\pjProjFixHi{9.8}
\def\pjPrompPsyW{92.5}
\def\pjPubTimeFranksW{91.8}
\def\pjPubTimeFranksF{74.2}
\def\pjPubTimePchipW{93.9}
\def\pjPubTimePchipF{89.1}
\def\pjExKpair{15}

% fig4m_text.tex: written by scripts/paper_v2_fig4m.py (repository); do not edit by hand.
\def\fmSThreeW{$+16.4$ ($+14.4$, $+18.3$)}
\def\fmSThreeF{$+7.9$ ($+6.1$, $+9.2$)}
\def\fmSThreeLevW{75.3}
\def\fmSThreeLevF{63.6}
\def\fmStartC{1.5}
\def\fmStartP{2.8}
\def\fmK{24}
\def\fmNclear{68}
\def\fmNall{75}
\def\fmNcell{70}
\def\fmNfix{68}
\def\fmNusers{200}
\def\fmNother{120}
\def\fmNtie{160}
\def\fmMainW{$+13.1$ ($+10.3$, $+14.7$)}
\def\fmMainF{$+8.8$ ($+6.8$, $+10.5$)}
\def\fmSingleW{$+2.6$ ($-0.5$, $+4.8$)}
\def\fmSingleF{$+3.4$ ($+0.5$, $+6.2$)}
\def\fmOtherW{$+16.0$ ($+12.6$, $+19.2$)}
\def\fmOtherF{$+12.0$ ($+8.5$, $+14.3$)}
\def\fmPeaks{$+8.1$}
\def\fmLaterW{$+12.1$}
\def\fmLaterF{$+10.7$}
\def\fmNoneF{$-2.0$ ($-5.3$, $+1.0$)}
\def\fmShapesF{$-3.1$ ($-4.8$, $-1.4$)}
\def\fmShapesW{$+13.7$ ($+10.8$, $+15.7$)}
\def\fmShapesNext{$+15.6$}
\def\fmOtherShapesW{$+5.0$}
\def\fmOtherShapesF{$+5.1$}
\def\fmPessW{$+12.6$ ($+10.0$, $+14.6$)}
\def\fmPessF{$-1.3$ ($-2.9$, $+0.3$)}
\def\fmDropSingle{17.2}
\def\fmDropPess{21.9}
\def\fmFixGap{8.6}
\def\fmBaldLo{1.9}
\def\fmBaldHi{5.7}


\renewcommand{\thesection}{S\arabic{section}}
\renewcommand{\theequation}{S\arabic{equation}}
\renewcommand{\thefigure}{S\arabic{figure}}
\renewcommand{\thetable}{S\arabic{table}}
\renewcommand{\thealgorithm}{S\arabic{algorithm}}
\algrenewcommand\algorithmiccomment[1]{\hfill{\footnotesize\color{icMuted}#1}}

\newcommand{\Nm}{\,Nm}
\newcommand{\z}{\mathbf z}
\newcommand{\code}[1]{\texttt{\small #1}}

\title{Supplementary Material\\[4pt]\large Accelerating Hip Exoskeleton Preference Optimization:\\
A Study with Simulated Users}
\author{Julian Seiler, Omid Mohseni, Niklas Sch\"afer, Mark Suppelt, Bastian Latsch, Mario Kupnik,\\
and Philipp Beckerle\\[2pt]\small julian.seiler@tu-darmstadt.de}
\date{}

\begin{document}
\maketitle

This supplement accompanies the paper above, submitted to ICORR 2027. It lists PsyPref step by step (Alg.~\ref{alg:psypref}),
one session with a simulated user (Alg.~\ref{alg:session}), the simulated users and every robustness condition with its
size and source (Tables~\ref{tab:user}--\ref{tab:conditions}), and the robustness results (Fig.~\ref{fig:robustness}).
Numbers without the prefix S refer to the paper. All results come from simulated users, and no person has used PsyPref.
The code, the simulation logs, and the instructions to reproduce every figure are in the repository
\url{https://github.com/TUDA-MUST/exo-psypref}, and each step below names the code that implements it.

\section{PsyPref Step by Step}
\label{sec:psypref}

PsyPref optimizes the $d = 8$ weights of a probabilistic movement primitive (ProMP)~\cite{paraschos}, a weighted sum of
basis functions over the gait cycle. It works in normalized coordinates $\z$, in which the prior over the weights is
standard normal. Every pair and the identified profile lie in the search region $\mathcal B = \{\z : \lVert\z\rVert \le
r\}$, where $r$ is the 99\,\% quantile of the norm of a standard normal draw in $d$ dimensions. The answer model reads a
choice of profile A over profile B as
\begin{equation}
P(\z_\mathrm{A}\succ\z_\mathrm{B}\mid\z^\ast,\ell,\lambda)=\frac{\lambda}{2}+(1-\lambda)\,\Phi\!\left(\frac{\Delta}{\ell^2}\right),
\qquad
\Delta=\tfrac12\bigl(\lVert\z_\mathrm{B}-\z^\ast\rVert_M^2-\lVert\z_\mathrm{A}-\z^\ast\rVert_M^2\bigr),
\label{eq:model}
\end{equation}
which is Eq.~(1) of the paper. Here $\z^\ast$ is the user's preferred profile, $\ell$ the discrimination scale, and
$\lambda$ the lapse rate~\cite{wichmann}, the probability that a user picks either profile at random. With a peak range, the metric is $M = J^\top J/101$, where $J$ is the Jacobian
of the 101-point torque curve with respect to $\z$ at $\z = \mathbf 0$, scaled so that its trace equals $d$. With fixed
peaks, $M = I$.

PsyPref holds its belief over $\mathbf x = (\z^\ast, \log\ell, \operatorname{logit}\lambda)$ as $N = 4000$ weighted
samples, called particles. Their priors are $\z^\ast\sim\mathcal N(\mathbf 0,I)$, $\log\ell\sim\mathcal N(0,0.5^2)$, and
$\lambda\sim\mathrm{Beta}(1,9)$. After each answer, every particle's weight is multiplied by the particle's likelihood of
that answer under Eq.~\eqref{eq:model}. When the effective sample size $1/\sum_i w_i^2$ falls below $N/2$, the particles
are resampled systematically and then moved by five random-walk Metropolis steps, which target the posterior given all
answers so far. The proposal covariance of these steps is $2.38^2/(d+2)$ times the particles' covariance.

Alg.~\ref{alg:psypref} lists one session. Most pairs come from a set of 220 candidates around the posterior mean
$\hat\z$, scored by how much the answer is expected to reduce the uncertainty about $\z^\ast$. The score is a Gaussian
proxy of the expected information gain~\cite{houlsby},
\begin{equation}
G = H - p\,H_\mathrm{A} - (1-p)\,H_\mathrm{B},\qquad H = \tfrac12\log\det\operatorname{Cov}(\z^\ast),
\label{eq:score}
\end{equation}
where $H$ is computed with the current weights, $H_\mathrm{A}$ and $H_\mathrm{B}$ with the weights after answer A or B,
and $p$ is the predicted probability of answer A, held within $[10^{-6}, 1-10^{-6}]$. The proxy ranks pairs and is not a
mutual information in nats. With probability 0.15, PsyPref instead asks a probe pair, one discrimination scale wide and
0.5, 1, or 2 scales from $\hat\z$, which mainly informs $\ell$ and $\lambda$.

Two operations keep pairs inside $\mathcal B$. $\mathrm{clip}(\mathbf c)$ moves a point outside $\mathcal B$ radially onto
its boundary. $\mathrm{trunc}(\mathbf c,\mathbf u,h)$ returns the ends $\mathbf c + t_+\mathbf u$ and $\mathbf c -
t_-\mathbf u$, each $t_\pm \le h$ as large as keeps that end inside $\mathcal B$. A pair therefore keeps its line and
direction and is never wider than $2h$. Candidates whose ends lie closer than $10^{-9}$ are dropped.

PsyPref delivers the posterior mean, clipped into $\mathcal B$. With a peak range, each delivered profile must assist both
extension and flexion. If one of its peaks falls below half of the range's lower limit, PsyPref delivers instead the
posterior draw closest to $\hat\z$, among 500 draws, that meets this limit. In the two peak ranges this replacement
occurred in 4 and 18\,\% of the sessions, never at a session's end. PsyPref needs about 0.15\,s per pair on one core.

In the main condition, simulated users always choose A or B. For the two conditions marked \textdaggerdbl{} in
Table~\ref{tab:conditions}, PsyPref treats ``no preference'' as an answer of its own, with a band of indecision around
$\Delta = 0$ whose half-width it learns from the first such answer on (\code{prompref/learning/ties.py}). Before that
answer, its session equals PsyPref's.

\begin{algorithm}[t]
\caption{One PsyPref session.}
\label{alg:psypref}
\begin{algorithmic}[1]
\Require prior over the weights, peak range or fixed peaks, number of comparisons $K$, seed
\State draw $N=4000$ particles $\mathbf x_i=(\z^\ast_i,\log\ell_i,\operatorname{logit}\lambda_i)$ from the prior;
  $w_i\gets 1/N$ \Comment{\code{smc.py}}
\For{$k=1,\dots,K$}
  \State $\hat\z\gets\sum_i w_i\z^\ast_i$;\quad $\hat\ell\gets\sum_i w_i\ell_i$
  \If{$U(0,1)<0.15$} \Comment{probe pair, \code{psychometric.py}}
    \State $\mathbf u\gets$ random unit vector;\quad $s\gets\pm1$;\quad $m\gets$ one of $\{0.5,1,2\}$
    \State $(\z_\mathrm{A},\z_\mathrm{B})\gets\mathrm{trunc}\bigl(\mathrm{clip}(\hat\z-s\,m\,\hat\ell\,\mathbf u),\,
      \mathbf u,\,\hat\ell/2\bigr)$
  \EndIf
  \If{no probe pair, or its ends coincide} \Comment{\code{candidates.py}}
    \State centers $\mathcal C\gets\{\hat\z\}\cup\{$3 particles drawn by weight$\}$, each clipped into $\mathcal B$
    \State directions $\mathcal U\gets\{d$ principal axes of the weighted covariance of $\z^\ast\}\cup\{$3 random unit
      vectors$\}$
    \State candidates $\gets\{\mathrm{trunc}(\mathbf c,\mathbf u,m\hat\ell/2):\mathbf c\in\mathcal C,\,
      \mathbf u\in\mathcal U,\,m\in\{4,2,1,0.5,0.25\}\}$ \Comment{$4\times11\times5=220$}
    \State $(\z_\mathrm{A},\z_\mathrm{B})\gets$ the candidate with the largest $G$, Eq.~\eqref{eq:score}
      \Comment{\code{acquisition.py}}
  \EndIf
  \State show the profiles of $\z_\mathrm{A}$ and $\z_\mathrm{B}$; receive the answer $y\in\{\mathrm{A},\mathrm{B}\}$
  \State $w_i\gets w_i\,P(y\mid\z_\mathrm{A},\z_\mathrm{B},\mathbf x_i)$, Eq.~\eqref{eq:model}; normalize
    \Comment{\code{smc.py}}
  \If{$1/\sum_i w_i^2<N/2$}
    \State resample systematically; five random-walk Metropolis steps on the posterior; $w_i\gets 1/N$
  \EndIf
\EndFor
\State $\z_\mathrm{out}\gets\mathrm{clip}(\hat\z)$ \Comment{\code{psychometric.py}}
\If{peak range and a peak of the profile of $\z_\mathrm{out}$ lies below half of the lower limit}
  \Comment{\code{amp\_free.py}}
  \State $\z_\mathrm{out}\gets$ the draw closest to $\hat\z$, among 500 drawn by weight, whose peaks both meet that limit
\EndIf
\State \Return the profile of $\z_\mathrm{out}$
\end{algorithmic}
\end{algorithm}

Alg.~\ref{alg:psypref} is implemented in \code{prompref/learning/} (\code{smc.py}, \code{acquisition.py},
\code{candidates.py}, \code{psychometric.py}) and \code{prompref/experiment/} (\code{stress.py} for the metric $M$,
\code{amp\_free.py} for $M$'s construction and the peak check). The library \code{psypref/session.py} wraps this code,
and with the data bundle its test reruns the paper's sessions to within $10^{-12}$\Nm.

\section{Simulated Users and Sessions}
\label{sec:users}

Each of the five individuals in the gait dataset~\cite{molinaro,molinarodata} provides 40 simulated users per peak-torque
setting, and the prior for these users comes from the other four individuals (Table~\ref{tab:user}). A user's preferred profile is
$\tau^\ast = a\,u^- + b\,u^+$, where $u^-$ and $u^+$ are the extension and flexion parts of the user's shape, each scaled
to a unit peak, and $a$ and $b$ are the preferred extension and flexion peaks. With a peak range, $a$ and $b$ are drawn
independently and log-uniformly within it. With fixed peaks, both equal 12\Nm.

A simulated user chooses between two profiles by their root-mean-square (RMS) distances $d$ to $\tau^\ast$, with noise:
\begin{equation}
P(\tau_\mathrm{A}\succ\tau_\mathrm{B})=\frac{\lambda}{2}+(1-\lambda)\,\Phi\!\left(\frac{d_\mathrm{B}-d_\mathrm{A}}{\sigma}\right),
\label{eq:user}
\end{equation}
which is Eq.~(2) of the paper, with $\sigma = 1$\Nm{} and $\lambda = 0.1$ in the main condition. Alg.~\ref{alg:session}
lists one session of such a user with any optimizer. One random stream per user, set by the user's answer seed, draws the
order in which the two profiles are shown and the answers. Every optimizer meets the same users with the same answer
seeds, so comparisons between optimizers are paired by user. The answers still differ between optimizers, because each
asks other pairs.

\begin{algorithm}[t]
\caption{One session of a simulated user with an optimizer, and its outcome.}
\label{alg:session}
\begin{algorithmic}[1]
\Require optimizer, preferred profile $\tau^\ast$, noise level $\sigma$, lapse rate $\lambda$, generic assistance
  $\tau_\mathrm g$, number of comparisons $K$, answer seed
\State random stream $\gets$ answer seed;\quad $d_0\gets\mathrm{RMS}(\tau_\mathrm g-\tau^\ast)$
  \Comment{\code{prompref/experiment/run.py}}
\For{$k=1,\dots,K$}
  \State $(\tau_\mathrm{A},\tau_\mathrm{B})\gets$ the optimizer's next pair
  \State draw which profile is shown first \Comment{matters only with a bias toward the first profile}
  \State $d_\mathrm{A}\gets\mathrm{RMS}(\tau_\mathrm{A}-\tau^\ast)$;\quad $d_\mathrm{B}\gets\mathrm{RMS}(\tau_\mathrm{B}-\tau^\ast)$
  \State answer A with probability $\lambda/2+(1-\lambda)\,\Phi\bigl((d_\mathrm{B}-d_\mathrm{A})/\sigma\bigr)$, else B
    \Comment{Eq.~\eqref{eq:user}, \code{simulated\_user.py}}
  \State give the answer to the optimizer
  \State $d_k\gets\mathrm{RMS}(\tau_k-\tau^\ast)$ for the profile $\tau_k$ the optimizer would deliver now
\EndFor
\State \Return $d_1,\dots,d_K$ and $P_\mathrm g(k)=\lambda/2+(1-\lambda)\,\Phi\bigl((d_0-d_k)/\sigma\bigr)$
\end{algorithmic}
\end{algorithm}

The outcome $P_\mathrm g(k)$ is the probability that the user chooses the profile identified after $k$ comparisons over
generic assistance, by the same rule~\eqref{eq:user}, and the paper averages it over users. Generic assistance is the
population's mean shape with both peaks at the geometric center of the peak range or at the fixed level.
$P_\mathrm g$ is 50\,\% for a profile as far from the preferred one as generic assistance, and at most 95\,\% at
$\lambda = 0.1$.

The user's just-noticeable difference (JND), the distance difference answered correctly in 75\,\% of attentive answers,
is $0.674\,\sigma$. No JND is published for hip assistance in walking. We therefore assume JNDs of
other joints and tasks and apply them to our users' preferred profiles to obtain $\sigma$ (Table~\ref{tab:anchors};
median RMS of these profiles 7.36 and 7.03\Nm{} at 3--35 and fixed 12\Nm). The main
condition lies near the static end, and the conditions reach $\sigma = 3$\Nm. Peng et al.\ report that their JND is
biased upward~\cite{peng}, so $\sigma = 3$\Nm{} is closer to an upper end than to an expected value.

\begin{table}[t]
\centering
\caption{The main simulated user.}
\label{tab:user}
\footnotesize
\begin{tabularx}{\textwidth}{@{}l>{\raggedright\arraybackslash}X>{\raggedright\arraybackslash}p{4.2cm}>{\raggedright\arraybackslash}p{3.4cm}@{}}
\toprule
Part & Model & Value & Source \\
\midrule
Users & Each fold holds out one of the five dataset individuals. The prior comes from the other four, the users from the
  held-out one & 5 folds $\times$ 40 users $=$ 200 users per setting & \cite{molinaro,molinarodata}, AB30--AB34 \\
Preferred profile & $\tau^\ast = a\,u^- + b\,u^+$, with $u^-$ and $u^+$ the unit extension and flexion parts of the
  user's shape & Peaks log-uniform in 3--12 or 3--35\Nm, or both at 12\Nm & Our choice \\
Personal shape & The held-out individual's hip moment, filtered at 10\,Hz, with a personal lag and personal shifts of
  its peaks and sign changes & Centre and spreads from the summary statistics of 11 participants, Student-t with
  $\nu = 10$ & \cite{schafer}, summary only \\
Distance & $d(c) = \mathrm{RMS}(c - \tau^\ast)$ over the gait cycle & -- & -- \\
Answer & Eq.~\eqref{eq:user} & $\sigma = 1$\Nm, $\lambda = 0.1$ & $\sigma$: Table~\ref{tab:anchors}. $\lambda$: our
  choice \\
JND & The distance difference answered correctly in 75\,\% of attentive answers & $0.674\,\sigma$ & Follows from the
  model \\
\bottomrule
\end{tabularx}
\end{table}

\begin{table}[t]
\centering
\caption{Anchors of the noise level $\sigma$: assumed JNDs applied to our users' preferred profiles.}
\label{tab:anchors}
\footnotesize
\begin{tabularx}{\textwidth}{@{}>{\raggedright\arraybackslash}Xp{4.6cm}cc@{}}
\toprule
Anchor & Source & Curve change (Nm RMS) & Implied $\sigma$ (Nm) \\
\midrule
Static force, 7\,\% & \cite{jones}\textdagger, \cite{pang}\textdagger & 0.49--0.52 & 0.73--0.77 \\
Static force, 10\,\% (our round-up) & None states 10\,\% & 0.71--0.74 & 1.05--1.09 \\
Walking stiffness, 25\,\% (a multiplier of about 3 from ankle stiffness) & \cite{maberry}\textdagger
  & 1.76--1.84 & 2.6--2.7 \\
Ankle timing, 2.8\,\%GC & \cite{peng} & 1.88--2.01 & 2.8--3.0 \\
\bottomrule
\multicolumn{4}{@{}l}{\textdagger\ Read as abstract only.}
\end{tabularx}
\end{table}

\FloatBarrier
\section{Robustness Conditions}
\label{sec:conditions}

Table~\ref{tab:conditions} lists every condition of Fig.~\ref{fig:robustness}. Each condition changes only what it names,
relative to the main user, and PsyPref keeps its model and prior throughout, apart from the ``no preference'' answers in
the two rows marked \textdaggerdbl. Hip data bear only on device noise and the spread of peak torques, so the other sizes
rest on other joints and tasks, or on our choice. Of the 42 values in Tables~\ref{tab:user}--\ref{tab:conditions}, 4 are
stated by a source, 1 lies inside a measured range, 14 are transferred from other joints or tasks, and 23 are our own
choices. The rows of our choice are stress tests, not estimates of real users. \code{CONDITIONS.md} in the repository
gives each condition's name in the code.

{\footnotesize
\begin{xltabular}{\textwidth}{@{}>{\raggedright\arraybackslash}p{3.0cm}>{\raggedright\arraybackslash}X
  >{\raggedright\arraybackslash}p{2.7cm}>{\raggedright\arraybackslash}p{1.6cm}>{\raggedright\arraybackslash}p{3.6cm}@{}}
\caption{Robustness conditions of Fig.~\ref{fig:robustness}, with model, size, class of the size, and source. Classes:
\emph{Stated}, a source gives this value for this quantity; \emph{Range}, inside a measured range; \emph{Transferred},
from another joint, task, or quantity; \emph{Ours}, no source gives the size.}
\label{tab:conditions}\\
\toprule
Label in Fig.~\ref{fig:robustness} & Model & Size & Class & Source and limits \\
\midrule
\endfirsthead
\toprule
Label in Fig.~\ref{fig:robustness} & Model & Size & Class & Source and limits \\
\midrule
\endhead
\bottomrule
\multicolumn{5}{@{}p{\textwidth}@{}}{\textdaggerdbl\ PsyPref uses ``no preference'' answers as data. These two rows use
users 8--39 of each fold, $n = 160$. \quad\textdagger\ Read as abstract only.}
\endlastfoot
\bottomrule
\endfoot
\multicolumn{5}{@{}l}{\textit{How they answer}} \\
Noise level $\sigma = 0.5$\Nm & $\sigma$ replaced & 0.5\Nm & Ours & Below every anchor in Table~\ref{tab:anchors} \\
Noise level $\sigma = 2$, 3\Nm & $\sigma$ replaced & 2, 3\Nm & Transferred & \cite{peng} (ankle timing),
  \cite{maberry}\textdagger{} (ankle stiffness), converted to Nm RMS by us \\
Lapse rate $\lambda = 0.2$, 0.3 & $\lambda$ replaced & 0.2, 0.3 & Transferred, ours & 0.2: indirect bound from validation
  rates in~\cite{arens,lee,alili} (other joints). 0.3: stress case above that bound \\
$\sigma$ differs per user (log SD 0.35) & $\sigma_u = \sigma\,e^{0.35 z}$, $z \sim \mathcal N(0,1)$ & Log SD 0.35 &
  Transferred & \cite{shepherd2018} (prosthesis stiffness, $n = 8$), \cite{miller}\textdagger{} (ankle
  timing): about fourfold range. \cite{shepherd2018} alone implies about 0.45 \\
Other answer-curve shapes & $\Phi$ replaced by a logistic, Student-t, or Cauchy curve, scaled to match the decisiveness
  of $\Phi$ on an earlier simulated population & Scales 0.5916, 0.8100, 0.4800; $\nu = 3$, $\nu = 1$ & Transferred
  (logistic), ours & Logistic: fitting convention in \cite{shepherd2018}. Heavy tails: stress cases \\
Device noise 0.6\Nm & Each presented curve gets its own Fourier series over harmonics 1--4, scaled to this RMS &
  0.6\Nm{} RMS & Stated & \cite{franks}, p.~e16-5, at the hips of a tethered emulator. The spectrum is ours \\
Bias toward first profile, 0.3\,$\sigma$ & $\beta\sigma$ added in favor of the profile shown first & $\beta = 0.3$ & Ours &
  No published estimate found \\
``No preference'' below 0.5\,JND\,\textdaggerdbl & The user answers ``no preference'' whenever $|d(A) - d(B)| <
  0.5$\,JND. A lapse is a coin flip between A and B & 0.5\,JND & Ours & \cite{arens} and~\cite{tucker} allow such
  answers, and neither gives a frequency \\
\midrule
\multicolumn{5}{@{}l}{\textit{How they change} (linear over the 40 comparisons of a session; familiarity holds from the
  first comparison)} \\
Fatigue: $\lambda$ from 0.05 to 0.3 & $\lambda$ rises & 0.05 to 0.30 & Ours & \cite{ingraham} rules out fatigue for their
  drift \\
$\sigma$ 1.5 to 0.75\Nm, and reverse & $\sigma$ falls or rises & 1.5 to 0.75\Nm, and reverse & Ours & Direction only:
  \cite{ingraham,elhusaini} \\
Preferred profile drifts 1, 2\,JND & $\tau^\ast$ grows by a factor, so that it ends this far from its start & 1, 2\,JND &
  Transferred & Direction and size: \cite{arens} (rise not significant, $P = 0.181$), \cite{ingraham}. 2\,JND exceeds the
  derived 1.0--1.6\,JND \\
Familiarity & The target becomes $(\tau^\ast + w\,\tau_h)/(1 + w)$, with $\tau_h = 0$, so it sits 20\,\% closer to zero
  torque & $w = 0.25$ & Ours & Direction only: \cite{diaz} ($n = 3$, ``possibly'') \\
\midrule
\multicolumn{5}{@{}l}{\textit{What they judge and prefer}} \\
Weber's law & $\sigma_\mathrm{eff} = \sigma\,(\mathrm{RMS}_\mathrm{pair} + a)/(\mathrm{RMS}_\mathrm{typ} + a)$ & $a = 1$\Nm,
  $\mathrm{RMS}_\mathrm{typ} = 4.83$\Nm & Transferred (form), ours (values) & \cite{hatzfeld} (fingertip vibration).
  $\mathrm{RMS}_\mathrm{typ}$ stems from an earlier population, so the median user here is also 1.38--1.43 times noisier \\
Over-assistance counts double & Squared error doubled where $|c| > |\tau^\ast|$, about 1.41 times in RMS & Factor 2 & Ours
  & No source. \cite{clites} points the other way \\
Judging single features & Distance over the two peak magnitudes and the two peak timings, $\sqrt{\sum_i
  (\Delta_i/\mathrm{JND}_i)^2}$, with magnitude JND $\max(f\,|\text{peak}|,\,0.5\text{\Nm})$, timing JND 2.8\,\%GC, and
  the scale matched to the main user's decisiveness & $f = 0.10$, 0.25 & Transferred & \cite{jones}\textdagger,
  \cite{pang}\textdagger{} state 7\,\%, static. \cite{maberry}\textdagger: walking about 3 times
  larger. \cite{peng} for timing. The 0.5\Nm{} floor is ours \\
Gait phases weighted unequally & Weighted RMS, $w(\varphi)$ with mean 1, on a per-user scale matched to the main user's
  decisiveness. Personal: $\log w$ a random Fourier series over harmonics 1--3, scaled to this SD over the cycle. Hip
  speed: $w \propto |\dot\theta| + 0.2\,\overline{|\dot\theta|}$ & Log SD 0.5, 1 & Ours & Rationale for the hip-speed
  form: power change $\Delta\tau\,\dot\theta$. No measurement of phase-dependent perception found \\
No single preferred profile & A random smooth function of a curve's first six harmonics, from random Fourier features,
  inside an envelope, on a matched per-user scale & Length scale 1.5, envelope 3 SD & Ours & Deliberately not in PsyPref's
  coordinates \\
Negative power disliked & $d(c) + \kappa\,\overline{\max(0, -c\,\dot\theta)}/\overline{|\dot\theta|}$ & $\kappa = 0$ with
  probability $1/2$, else log-uniform in 0.5--2 & Ours & Direction only: \cite{ramella} \\
\midrule
\multicolumn{5}{@{}l}{\textit{How they differ from the prior}} \\
Users from other datasets & Shapes from other individuals, 40 users per group & Young and older adults on a treadmill at
  comfortable speed; level and inclined walking at 1.0\,m/s & Stated & \cite{fukuchi,reznick}. These are biological
  moments, and treating them as preferred shapes is our assumption \\
Peaks unlike the prior & Peak laws: three narrow cohorts, flexion or extension larger, correlated & Medians at 0.2, 0.5,
  0.8 of the log range with CV 0.2; ratio 1.5; correlation 0.8 & Range (CV), ours & CV: \cite{ramella} (0.13--0.15),
  \cite{arens} (0.30--0.34), \cite{shepherd2018} (0.30). Rest is ours \\
Flexion peak 6, 12\,\%GC later & Piecewise-linear phase warp that moves the flexion peak and keeps the extension peak &
  6, 12\,\%GC & Transferred & \cite{kim2022} (soft flexion exosuit, metabolic optimum, $n = 8$) \\
Other shapes unlike the prior & Surrogates: the preferred curve's magnitudes kept, its phases randomized, the result
  clipped to the setting's top & Harmonics 1--50 & Transferred & \cite{theiler}, a method for testing time series. It
  serves here as a source of curves without the prior's shape \\
 & Heavy tails: personal lag and shifts with $\nu = 1.5$ instead of 10, same scales & $\nu = 1.5$ & Ours & Stress case:
  infinite variance before invalid draws are rejected \\
 & Family mixture: a share of users replaced, half from Sch\"afer et al.'s family, half from Franks et al.'s & 20\,\% &
  Ours & -- \\
Shapes from the 4--6-param.\ family & Four timings from a Student-t distribution in Sch\"afer et al.'s family. Draws
  outside the device limits are rejected & Mean and diagonal covariance of 11 participants, $\nu = 10$ & Stated (family),
  ours ($\nu$) & \cite{schafer}. CMA-ES, 4--6 param., searches this family \\
\midrule
Pessimistic user (six at once)\,\textdaggerdbl & Six of the above at once & $\sigma = 2$\Nm, fatigue, bias 0.3\,$\sigma$, no
  preference at 0.5\,JND, phase weight log SD 1, Sch\"afer-family shapes & Ours & Combination of the rows above \\
\end{xltabular}
}

\begin{figure}[t]
\centering
\scalebox{1.3}{% fig4m.tex: Fig. 4, variant 4M (4L as a measurement engineer would read it: saturation marked, the known
% systematic part shown, intervals that hold for all cells at once). One column, 252 pt wide.
%   \input{figures/tikz_v2/fig4m.tex}   inside \begin{figure}[t]\centering ... \end{figure}
% Preamble: \usepackage{pgfplots}\pgfplotsset{compat=1.18}\input{icorr_style.tex}
% Data: \icorrtikzdata (default figures/tikz_v2/data/), written by scripts/paper_v2_fig4m.py (amp-free
% repository): fig4m_rows.csv (one line per group header or row: y in row units; per setting the light and
% the blue five values, nan where the cell has the other colour or does not exist; the cap marker; the
% label), fig4m_meta.csv (axis limits and the layout in pt).
% Forest plot: PsyPref minus the strongest optimizer of that row, picked again in every bootstrap draw,
% paired by user, chosen over generic assistance after 24 comparisons, in points, each user rated with their
% own noise level and lapse rate at the session's end. Circle 3-35 Nm (0.25 rows above the row's centre),
% square fixed 12 Nm (0.25 below). Thick bar: 95 % interval of that row and setting. Thin bar: 95 % for the 73
% cells jointly. Blue with its value where the thin bar reaches 0, else icBlueLight. Open circle with "both
% at 95 %": PsyPref and the comparator at the measure's maximum in every draw. "(n)": the smallest lead of n
% conditions per setting; dagger: PsyPref with tie handling (n = 160). The first and the last row are
% bold, no header; the second row is the main condition against Fig. 2's comparator. Row labels give a
% deviation's size where one number suffices.
\providecommand{\icorrtikzdata}{figures/tikz_v2/data/}%
\makeatletter
\providecommand{\icorrload}[2]{% \icorrload{name}{file.csv}: every cell as \icv{name}{row}{column}%
  \pgfplotstableread[col sep=comma]{\icorrtikzdata #2}\ic@table
  \pgfplotstableforeachcolumn\ic@table\as\ic@col{%
    \pgfplotstableforeachcolumnelement{\ic@col}\of\ic@table\as\ic@cell{%
      \expandafter\global\expandafter\let\csname ic@#1@\pgfplotstablerow @\ic@col\endcsname
        \ic@cell}}}%
\providecommand{\icv}[3]{\csname ic@#1@#2@#3\endcsname}%
\def\fmlabel#1#2{\expandafter\def\csname fm@lab@#1\endcsname{#2}}%
\def\fmlab#1{\csname fm@lab@#1\endcsname}%
\makeatother
% the row labels, by key (fig4m_rows.csv's column key); condition values from fig4m_meta.csv
\fmlabel{MAIN}{Main condition (Fig.~2)}%
\fmlabel{G1}{How they answer}%
\fmlabel{R1}{Noise level $\sigma$ = \pgfmathprintnumber{\icv{fmm}{0}{s_r1}}\,Nm}%
\fmlabel{R2}{Noise level $\sigma$ = \pgfmathprintnumber{\icv{fmm}{0}{s_r2}}, \pgfmathprintnumber{\icv{fmm}{0}{s_r3}}\,Nm}%
\fmlabel{AL}{Lapse rate $\lambda$ = \pgfmathprintnumber{\icv{fmm}{0}{lapse_a2}}, \pgfmathprintnumber{\icv{fmm}{0}{lapse_r3}}}%
\fmlabel{A1}{$\sigma$ differs per user (log SD \icv{fmm}{0}{a1_logsd})}%
\fmlabel{A6}{Other answer-curve shapes}%
\fmlabel{A3}{Device noise \icv{fmm}{0}{a3_noise}\,Nm}%
\fmlabel{A4}{Bias toward first profile, \icv{fmm}{0}{a4_bias}\,$\sigma$}%
\fmlabel{A5T}{``No preference'' below \icv{fmm}{0}{a5_nopref}\,JND}%
\fmlabel{G2}{How they change}%
\fmlabel{A7}{Fatigue: $\lambda$ from \icv{fmm}{0}{a7_from} to \icv{fmm}{0}{a7_to}}%
\fmlabel{A8}{$\sigma$ \icv{fmm}{0}{a8_from} to \icv{fmm}{0}{a8_to}\,Nm, and reverse}%
\fmlabel{B4}{Preferred profile drifts \icv{fmm}{0}{b4_a}, \icv{fmm}{0}{b4_b}\,JND}%
\fmlabel{B3}{Familiarity}%
\fmlabel{G3}{What they judge and prefer}%
\fmlabel{B1}{Weber's law}%
\fmlabel{B2}{Over-assistance counts double}%
\fmlabel{B5}{Judging single features}%
\fmlabel{B67}{Gait phases weighted unequally}%
\fmlabel{B8}{No single preferred profile}%
\fmlabel{R4}{Negative power disliked}%
\fmlabel{G4}{How they differ from the prior}%
\fmlabel{R6}{Users from other datasets}%
\fmlabel{PM}{Peaks unlike the prior}%
\fmlabel{PS}{Flexion peak 6, 12\,\%GC later}%
\fmlabel{SO}{Other shapes unlike the prior}%
\fmlabel{SC}{Shapes from the 4--6-param. family}%
\fmlabel{C4T}{Pessimistic user (\icv{fmm}{0}{ndev_word} at once)}%
\begin{tikzpicture}[x=1pt, y=1pt]
\pgfkeys{/pgf/number format/assume math mode=true}% numbers in the text font (Times)
\icorrload{fmm}{fig4m_meta.csv}%
\icorrload{fmr}{fig4m_rows.csv}%
\pgfmathtruncatemacro\fmN{\icv{fmm}{0}{n_rows}-1}%
\pgfmathsetmacro\fmL{\icv{fmm}{0}{lab_w}}\pgfmathsetmacro\fmR{\icv{fmm}{0}{ax_r}}%
\pgfmathsetmacro\fmP{0.85*\icv{fmm}{0}{pitch}}\pgfmathsetmacro\fmB{\icv{fmm}{0}{bottom}}%
\pgfmathsetmacro\fmoff{\icv{fmm}{0}{off}}\pgfmathsetmacro\fmmr{\icv{fmm}{0}{mark_r}}%
\pgfmathsetmacro\fmms{\icv{fmm}{0}{mark_s}}\pgfmathsetmacro\fmvgapx{\icv{fmm}{0}{val_gapx}}%
\pgfmathsetmacro\fmxmin{\icv{fmm}{0}{xmin}}\pgfmathsetmacro\fmxmax{\icv{fmm}{0}{xmax}}%
\pgfmathsetmacro\fmymin{\icv{fmm}{0}{ymin}}\pgfmathsetmacro\fmymax{\icv{fmm}{0}{ymax}}%
\pgfmathsetmacro\fmscale{\icv{fmm}{0}{scale}}%
\pgfmathsetmacro\fmH{(\fmymax-\fmymin)*\fmP}\pgfmathsetmacro\fmW{\fmR-\fmL}%
\pgfmathsetmacro\fmylab{\icv{fmm}{0}{ylabels}}\pgfmathsetmacro\fmlinetop{\icv{fmm}{0}{linetop}}%
\def\fmlight{icBlueLight}% a lead whose thin bar lies above 0: the lighter PsyPref
\pgfplotsset{
  fm mark w/.style={mark=*, mark size=\fmmr pt},
  fm mark f/.style={mark=square*, mark size=\fmms pt},
  fm marks/.style={only marks, mark options={fill=#1, draw=#1, line width=0.3pt}},
  % error bars only: no line through the points and no marks
  fm bar/.style 2 args={draw=none, mark=none, error bars/.cd, x dir=both, x explicit, error mark=none,
    error bar style={draw=#1, line width=#2, line cap=butt}, /pgfplots/.cd},
}
\begin{axis}[ic axis, scale only axis, at={(\fmL pt,\fmB pt)}, anchor=south west,
  width=\fmW pt, height=\fmH pt,
  xmin=\fmxmin, xmax=\fmxmax, ymin=\fmymin, ymax=\fmymax, y dir=reverse,
  axis y line=none, ytick=\empty,
  xtick={-10,0,10,20,30},
  xticklabels={\textminus10, 0, +10, +20, +30},
  xticklabel style={inner sep=1pt},
]
% no lead (0): the strongest published optimizer
\draw[icMuted, line width=0.5pt] (axis cs:0,\fmlinetop) -- (axis cs:0,\fmymax);
\node[anchor=base west, inner sep=0pt, font=\scriptsize, text=icInk, xshift=-2.5pt]
  at (axis cs:0,\fmylab) {Favors PsyPref\,$\rightarrow$};
% thin bars, thick bars, then markers: light first, blue on top; circle 3-35, square fixed 12
\foreach \c/\t/\s/\m in {\fmlight/in/w/-, \fmlight/in/f/+, icBlue/out/w/-, icBlue/out/f/+}{%
  \edef\fmplot{\noexpand\addplot[fm bar={\c}{0.5pt}, forget plot] table[col sep=comma, x=\t_\s,
    y expr=\noexpand\thisrow{y}\m\fmoff,
    x error plus expr=\noexpand\thisrow{\t_jhi_\s}-\noexpand\thisrow{\t_\s},
    x error minus expr=\noexpand\thisrow{\t_\s}-\noexpand\thisrow{\t_jlo_\s}]
    {\icorrtikzdata fig4m_rows.csv};}\fmplot
  \edef\fmplot{\noexpand\addplot[fm bar={\c}{1pt}, forget plot] table[col sep=comma, x=\t_\s,
    y expr=\noexpand\thisrow{y}\m\fmoff,
    x error plus expr=\noexpand\thisrow{\t_hi_\s}-\noexpand\thisrow{\t_\s},
    x error minus expr=\noexpand\thisrow{\t_\s}-\noexpand\thisrow{\t_lo_\s}]
    {\icorrtikzdata fig4m_rows.csv};}\fmplot
  \edef\fmplot{\noexpand\addplot[fm marks=\c, fm mark \s, forget plot] table[col sep=comma, x=\t_\s,
    y expr=\noexpand\thisrow{y}\m\fmoff] {\icorrtikzdata fig4m_rows.csv};}\fmplot}
% both methods at the cap: an open circle, nothing measured
\addplot[only marks, fm mark w, mark options={fill=white, draw=icBlue, line width=0.6pt}, forget plot]
  table[col sep=comma, x=cap_w, y expr=\thisrow{y}-\fmoff] {\icorrtikzdata fig4m_rows.csv};
\end{axis}
% labels of the blue cells: the value left of the thin bar, "both at 95 %" right of its circle, centred on their
% own marker's line, shifted by the script's vdy (pt, up positive) where something drawn is close
\def\fmvalue#1#2#3{% row, setting tag, marker offset (row units)
  \pgfmathsetmacro\fmvy{\fmB+(\fmymax-\icv{fmr}{#1}{y}-(#3))*\fmP+\icv{fmr}{#1}{vdy_#2}}%
  \ifnum\icv{fmr}{#1}{pos_#2}=3
    \node[inner sep=0pt, font=\scriptsize, text=icBlue, anchor=east] at (\icv{fmr}{#1}{vx_#2},\fmvy)
      {\ifnum\icv{fmr}{#1}{neg_#2}=1\textminus\fi\icv{fmr}{#1}{lab_#2}};%
  \fi
  \ifnum\icv{fmr}{#1}{pos_#2}=4
    \node[inner sep=0pt, font=\scriptsize, text=icBlue, anchor=west] at (\icv{fmr}{#1}{vx_#2},\fmvy) {both at 95\,\%};%
  \fi}%
% group headers and row labels, left of the axis
\foreach \i in {0,...,\fmN}{%
  \pgfmathsetmacro\fmy{\fmB+(\fmymax-\icv{fmr}{\i}{y})*\fmP}%
  \ifnum\icv{fmr}{\i}{is_header}=1
    \node[anchor=mid west, inner sep=0pt, font=\scriptsize\bfseries, text=icInk]
      at (0,\fmy) {\fmlab{\icv{fmr}{\i}{key}}};%
  \else
    \ifnum\icv{fmr}{\i}{bold}=1 \def\fmfont{\scriptsize\bfseries}\def\fmx{0}%
    \else \def\fmfont{\scriptsize}\def\fmx{5}\fi
    \node[anchor=mid west, inner sep=0pt, font=\fmfont, text=icInk]
      at (\fmx,\fmy) {\fmlab{\icv{fmr}{\i}{key}}%
        \ifnum\icv{fmr}{\i}{n_merged}>1\ (\icv{fmr}{\i}{n_merged})\fi
        \ifnum\icv{fmr}{\i}{tie}=1\,\textdagger\fi};%
    \fmvalue{\i}{w}{-\fmoff}%
    \fmvalue{\i}{f}{\fmoff}%
  \fi}
% key (shape = the setting's peak torque), on the line of the zero line's label, left
\pgfmathsetmacro\fmyk{\fmB+(\fmymax-\fmylab)*\fmP}%
\begin{scope}[font=\scriptsize, every node/.style={anchor=base west, inner sep=0pt, text=icInk}]
  \node (fmk1) at (0,\fmyk) {Peak torque};
  \fill[\fmlight] ([xshift=4.4pt, yshift=1.7pt]fmk1.base east) circle[radius=\fmmr pt];
  \node (fmk2) at ([xshift=7.3pt]fmk1.base east) {3--35\,Nm};
  \fill[\fmlight] ([xshift=4.4pt-\fmms pt, yshift=1.7pt-\fmms pt]fmk2.base east)
    rectangle ([xshift=4.4pt+\fmms pt, yshift=1.7pt+\fmms pt]fmk2.base east);
  \node at ([xshift=7.3pt]fmk2.base east) {Fixed 12\,Nm};
\end{scope}
% x label, one line, centred under the zero line
\pgfmathsetmacro\fmzero{\fmL-\fmxmin*\fmscale}%
\node[anchor=base, inner sep=0pt, font=\footnotesize, text=icInk] at (\fmzero,1.9)
  {$\Delta P_\mathrm{g}$ (percentage points)};
\pgfresetboundingbox
\useasboundingbox (0,0) rectangle (252,\fmB+\fmH);
\end{tikzpicture}%
}
\caption{Difference in $P_\mathrm{g}$ between PsyPref and the best published optimizer after \fmK\ comparisons, per
violated assumption, in simulation. Positive values favor PsyPref. Markers: paired difference from the best of Arens'
GP~\cite{arens} (also from the prior's mean, no gate), CMA-ES~\cite{sortcma} (also with population 4),
LineCoSpar~\cite{linecospar}, and \mbox{CMA-ES, 4--6~param.}, picked per bootstrap draw. Thick bars: 95\,\% per row. Thin
bars: 95\,\% across all \fmNall\ cases (condition $\times$ peak-torque setting)~\cite{besag}, dark blue where they reach
0. Open circle: both at $P_\mathrm{g}$'s ceiling of 95\,\%. (2)--(6): smallest difference of that many conditions. \textdagger: with the
indecision band. Table~\ref{tab:conditions} lists every condition with its model, size, and source. The paper's Fig.~4
summarizes this figure.}
\label{fig:robustness}
\end{figure}

\bibliographystyle{IEEEtran}
\bibliography{refs}

\end{document}
